\documentclass[10pt,a4paper]{amsart}
\usepackage[margin=19mm]{geometry}
\usepackage{amsmath,amssymb,mathtools}
\usepackage[scr=rsfs,cal=euler]{mathalfa}
\usepackage{xfrac}
\usepackage[unicode,hidelinks,bookmarks=false]{hyperref}
\newcommand{\ie}{i.e.\ }
\newcommand{\cf}{cf.\ }
\newcommand{\resp}{respectively\ }
\newcommand{\Subsection}{\subsection}
\providecommand{\nobreakdash}{\nobreak\hbox{-}}
\setlength{\emergencystretch}{2em}
\multlinegap0pt
\usepackage{bm}
\usepackage{bbm}
\usepackage{dsfont}
\usepackage{xspace}
\usepackage{cancel}
\usepackage{mathtools}
\usepackage{amsthm}
\makeatletter
\@ifundefined{cor}{\newtheorem{cor}{Cor}}{}
\@ifundefined{lemma}{\newtheorem{lemma}[cor]{Lemma}}{}
\makeatother
\def\funch{h}
\newcommand{\R}{\mathbb{R}}
\title[A surface finite element scheme for a stochastic PDE on an e]{A surface finite element scheme for a stochastic PDE on an evolving curve}
\author{Paola Pozzi, Bj\"orn Stinner}
\date{Source: arXiv:2507.01527v1 (CC BY 4.0)}
\begin{document}
\maketitle

\noindent\emph{Source and licence.} A surface finite element scheme for a stochastic PDE on an evolving curve. Version arXiv:2507.01527v1 (CC BY 4.0), distributed under the Creative Commons Attribution 4.0 International licence, as recorded in the arXiv licence metadata for this exact version and shown on the abstract page https://arxiv.org/abs/2507.01527v1. Full rights notice: licenses/arxiv-2507.01527-CC-BY-4.0.txt.\par\smallskip

\noindent\textit{Editorial scope.} Counted unit: the Lemma whose source label is \texttt{RH1} in the article (source0.tex lines 627--651, source label \texttt{RH1}), delivered together with the article's own complete original proof of it (source0.tex lines 653--786, one \texttt{proof} environment of 134 lines). No mathematics is written, repaired or re-derived: statement and proof are the article's own text line for line. Required context retained uncounted: regularity (lines 163--165); (4.2) (lines 426--433); (4.2)bis (lines 426--433); (4.2)extra (lines 426--433); Ritz1 (lines 617--621); Ritz2 (lines 617--621). Every definition, display and label of those lines is retained unchanged. Omitted from this package and not redistributed: 1--162, 166--425, 434--616, 622--626, 652--652, 787--1495. Nothing inside a retained excerpt is omitted or altered: every retained body is delivered exactly as the article's lines are written, so no omission receipt of any kind accompanies this package. Citations: the retained excerpts carry no citation command at all, so no bibliography entry of the article is transcribed and no reference is redistributed. Reference adaptation: none was needed and none was made; every reference inside the delivered unit resolves inside it, so no unresolved reference or citation is produced. Printed slips kept literal: no printed word, symbol, hypothesis or number of the counted statement or of its proof was altered. Layout and class: the article's own class is \texttt{article} with packages including babel, amssymb, amsmath, indentfirst, graphicx, epsfig, lscape, hyperref, color, soul, xcolor, amsthm; the wrapper is the lane's pinned \texttt{amsart} editorial wrapper at 10pt on A4, which restates the theorem-like environments the retained text needs and adds the article's own macro definitions that the retained text needs (\texttt{\textbackslash R}, \texttt{\textbackslash funch}), with extra wrapper packages (bm, bbm, dsfont, xspace, cancel, mathtools). The delivered numbering is the unit's own. Assets: no retained span contains a figure environment, a graphics-inclusion command, a PDF-page inclusion, a table or a TikZ picture; no image file is copied into this package, the package declares no asset dependency and no image file is redistributed with it. Licence: version arXiv:2507.01527v1 of this article is distributed under the Creative Commons Attribution 4.0 International licence, as recorded in the arXiv licence metadata for this exact version and shown on the abstract page https://arxiv.org/abs/2507.01527v1. A full rights notice is delivered at licenses/arxiv-2507.01527-CC-BY-4.0.txt. Characters: every retained excerpt is pure ASCII, so the delivered engine, XeLaTeX with its default Latin Modern text font, reports no missing character.
\begin{align}\label{regularity}
 0<d_{0} \leq |u_{x}(t,x)| \leq \frac{1}{d_{0}} \qquad \text{ uniformly on } [0,T] \times S^{1}
\end{align}

\begin{align}
\label{(4.2)}
\| v-I_h v \|_{L^2(S^1)} &\leq C h^k \|v \|_{W^{k,2}(S^1)} &\text{ for $k=1,2$}\,,\\
\label{(4.2)bis}
\| (v-I_h v)_x \|_{L^2 (S^1)} & \leq C h \| v \|_{W^{2,2} (S^1)}\,,\\
\label{(4.2)extra}
 \| (I_{h} v)_{x} \|_{L^{2} (S^{1})} & \leq C \| v_{x} \|_{ L^{2}(S^{1})}\, . 
\end{align} 

\begin{align} \label{Ritz1}
 \| (z - \mathcal{R}_{h} (z) )_{x}\|_{H} \leq C \| \partial_{s} z - \partial_{s} \mathcal{R}_{h} (z) \|_{L^{2}(ds(t))} &\leq C h^{s-1} \|z\|_{W^{s,2}(S^{1})} \qquad s=1,2,\\
\label{Ritz2}
 \| z - \mathcal{ R}_{h} (z) \|_{H} \leq C \| z - \mathcal{R}_{h} (z) \|_{L^{2}(ds(t))} &\leq C h^{s} \|z\|_{W^{s,2}(S^{1})} \qquad s=1,2.
\end{align}

\begin{lemma}[Properties of $\mathcal{R}_{h}$]
Let $z(t) \in W^{s,2}(S^{1})$ for any $t \in [0,T]$ and for $s \in \{ 1,2 \}$. Then, for any $t, \tau \in [0,T]$ we have
\begin{align}
\label{RH1}
 \| (\mathcal{R}_{h} (z(t)) -\mathcal{R}_{h} (z(\tau)) )_{x} \|_{H}& \leq C_{g} |t -\tau| h^{s-1}\| z(\tau)\|_{W^{s,2}(S^{1})} + C \| (z(t) -z(\tau) )_{x} \|_{H},\\
 \label{RH2}
 \| \mathcal{R}_{h} (z(t)) -\mathcal{R}_{h} (z(\tau)) \|_{H}& \leq C_{g} |t -\tau| \| z(\tau)\|_{V}+
 C\| z(t) -z(\tau) \|_{V}.
\end{align}
More precisely, we have that
\begin{multline}\label{RH3}
\| (\mathcal{R}_{h} (z(t)) -\mathcal{R}_{h} (z(\tau)) - (z(t) -z(\tau)) )_{x}\|_{H} \\
\leq Ch^{s-1}
\| z(t) -z(\tau)\|_{W^{s,2}(S^{1})} + C_{g} |t -\tau|h^{s-1}
\| z(\tau)\|_{W^{s,2}(S^{1})} ,
\end{multline}
and 
\begin{multline}\label{RH4}
\| \mathcal{R}_{h} (z(t)) -\mathcal{R}_{h} (z(\tau)) - (z(t) -z(\tau)) \|_{H} \\
 \leq C_{g}|\tau-t|h^{s}\| z(\tau)\|_{W^{s,2}(S^{1})} + Ch^{s}
\| z(t) -z(\tau)\|_{W^{s,2}(S^{1})} .
\end{multline}
The constant $C_{g}$ depends on the evolution of the geometry.
In particular, if the velocity $u_{x}$ does not change in time (i.e. the curve is stationary) then $C_{g}=0$.
\end{lemma}

\begin{proof}
By definition for any time $t \in [0,T]$ and $\varphi_{h} \in S_{h}$ we have that
$$\int_{S^{1}} (z(t))_{x} \frac{\varphi_{hx}}{|u_{x}(t)|} dx = \int_{S^{1}} ( \mathcal{R}_{h} z(t))_{x} \frac{\varphi_{hx}}{|u_{x}(t)|} dx $$
therefore, for any $t, \tau \in [0,T]$ we have
\begin{align}\label{start}
\int_{S^{1}}& (\mathcal{R}_{h} (z(t)) -\mathcal{R}_{h} (z(\tau)) )_{x} \frac{\varphi_{hx}}{|u_{x}(t)|}dx \\
&= \int_{S^{1}} (z(t) -\mathcal{R}_{h} (z(\tau)) )_{x} \varphi_{hx} \left(\frac{1}{|u_{x}(t)|} - \frac{1}{|u_{x}(\tau)|}\right)dx
+ \int_{S^{1}} (z(t) -z(\tau) )_{x} \frac{\varphi_{hx}}{|u_{x}(\tau)|}dx . \notag
\end{align}
Choosing $\varphi_{h}=\mathcal{R}_{h} (z(t)) -\mathcal{R}_{h} (z(\tau)) $ and using \eqref{regularity} and the regularity of $u$ we obtain
\begin{align*}
\| (\mathcal{R}_{h} (z(t)) -\mathcal{R}_{h} (z(\tau)) )_{x} \|_{H}^{2} &\leq
\int_{S^{1}} |(\mathcal{R}_{h} (z(t)) -\mathcal{R}_{h} (z(\tau)) )_{x}|^{2} \frac{1}{|u_{x}(t)|}dx \\
&\leq C |t -\tau|\| ( z(t) -z(\tau) + z(\tau) -\mathcal{R}_{h} (z(\tau)) )_{x} \|_{H} \| (\mathcal{R}_{h} (z(t)) -\mathcal{R}_{h} (z(\tau)) )_{x} \|_{H}\\
& \quad + C \| (z(t) -z(\tau) )_{x} \|_{H} \| (\mathcal{R}_{h} (z(t)) -\mathcal{R}_{h} (z(\tau)) )_{x} \|_{H}.
\end{align*}
Inequality \eqref{RH1} follows from \eqref{Ritz1}.

Next, we use a Poincar\'{e}-inequality for a map $\funch:S^{1} \to \R$ along the curve $u(t)$:
$$\| \funch - \bar{h}\|_{L^{2}(ds(t))} \leq C\| \partial_{s} \funch \|_{L^{2}(ds(t))} \qquad \text{ where } \bar{h}:= \frac{1}{|S^{1}|}\int_{S^{1}} \funch ds(t)= \frac{1}{|S^{1}|}\int_{S^{1}} \funch |u_{x}(t)| dx.$$
Note that $C$ depends on the length of the curve $u(t)$, which is uniformly bounded in $t$ by \eqref{regularity}. We obtain, again with the help of \eqref{regularity}, that
\begin{align*}
\|\funch\|_{H} \leq C \| \funch\|_{L^{2}(ds(t))} \leq C\| \funch - \bar{h}\|_{L^{2}(ds(t))} + C\| \bar{h}\|_{L^{2}(ds(t))} \leq C \| \funch_{x}\|_{H} + C |\bar{h}|.
\end{align*}
Choosing $\funch:=\mathcal{R}_{h} (z(t)) -\mathcal{R}_{h} (z(\tau)) $ and using \eqref{RH1} we obtain \eqref{RH2} since
\begin{align*}
|\bar{h}| &\leq C \Big|\int_{S^{1}} \mathcal{R}_{h} (z(t)) |u_{x}(t)| - \mathcal{R}_{h} (z(\tau)) |u_{x}(t)| dx \Big|\\
&\leq C \Big|\int_{S^{1}} \mathcal{R}_{h} (z(t)) |u_{x}(t)| - \mathcal{R}_{h} (z(\tau)) |u_{x}(\tau)| dx \Big| + C \Big|\int_{S^{1}} \mathcal{R}_{h} (z(\tau)) ( |u_{x}(\tau)| - |u_{x}(t)|) dx \Big|\\
& \leq C \Big|\int_{S^{1}} z(t) |u_{x}(t)| - z(\tau) |u_{x}(\tau)| dx \Big| + C |t -\tau| \| \mathcal{R}_{h} (z(\tau)) \|_{L^{1}(S^{1})}\\
& \leq C \| z(t) -z(\tau)\|_{H} + C |t -\tau| ( \| z(\tau)\|_{H} +\| \mathcal{R}_{h} (z(\tau)) \|_{H} )\\
&\leq C \| z(t) -z(\tau)\|_{H} + C |t -\tau| ( \| z(\tau)\|_{H} +\| \mathcal{R}_{h} (z(\tau)) -z(\tau)\|_{H} )\\
& \leq C \| z(t) -z(\tau)\|_{H} + C |t -\tau| (\| z(\tau)\|_{H} +h^{s}\| z(\tau)\|_{W^{s,2}(S^{1})})
\end{align*}
where we have used \eqref{Ritz2} in the last inequality.

%%%%%%%%%
Next, starting again from \eqref{start} we obtain for $\varphi_{h} \in S_{h}$ 
\begin{align*}
\int_{S^{1}}& [(\mathcal{R}_{h} (z(t)) -\mathcal{R}_{h} (z(\tau))) -(z(t) -z(\tau)) ]_{x} \frac{\varphi_{hx} }{|u_{x}(t)|}dx \\
&= \int_{S^{1}} (z(t) -\mathcal{R}_{h} (z(\tau)) )_{x} \varphi_{hx} \left(\frac{1}{|u_{x}(t)|} - \frac{1}{|u_{x}(\tau)|}\right)dx \\
& \quad
+ \int_{S^{1}} (z(t) -z(\tau) )_{x} \varphi_{hx} \left( \frac{1}{|u_{x}(\tau)|} - \frac{1}{|u_{x}(t)|} \right)dx .
\end{align*}
Choosing $\varphi_{h}=(\mathcal{R}_{h} (z(t)) -\mathcal{R}_{h} (z(\tau))) -I_{h}(z(t) -z(\tau))=
(\mathcal{R}_{h} (z(t)) -\mathcal{R}_{h} (z(\tau))) - (z(t) -z(\tau))+(z(t) -z(\tau)) -I_{h}(z(t) -z(\tau))$
we obtain, setting for simplicity of exposition
\begin{align*}
R_{h}:=\mathcal{R}_{h} (z(t)) -\mathcal{R}_{h} (z(\tau)), \qquad Z:=z(t) -z(\tau), \qquad I_{h}Z:= I_{h}(z(t) -z(\tau))=I_{h}z(t) -I_{h}z(\tau),
\end{align*}
and using \eqref{regularity}
\begin{align*}
d_{0} \|(R_{h}-Z)_{x}\|_{H}^{2}
&\leq
\int_{S^{1}}\frac{|(R_{h}-Z)_{x}|^{2}}{|u_{x}(t)|} dx=\int_{S^{1}}\frac{(R_{h}-Z)_{x} (I_{h}Z-Z)_{x}}{|u_{x}(t)|}\\
& \quad 
 + \int_{S^{1}} (z(\tau) -\mathcal{R}_{h} (z(\tau)) )_{x} (R_{h}-I_{h}Z)_{x}\left(\frac{1}{|u_{x}(t)|} - \frac{1}{|u_{x}(\tau)|}\right)dx \\
& \leq \epsilon \|(R_{h}-Z)_{x}\|_{H}^{2} + C_{\epsilon} \|(I_{h}Z-Z)_{x}\|_{H}^{2} \\
& \quad + C |t -\tau|
\| (z(\tau) -\mathcal{R}_{h}(z(\tau)))_{x}\|_{H} ( \|(R_{h}-Z)_{x}\|_{H} + \|(I_{h}Z-Z)_{x}\|_{H})\\
& \leq \epsilon \|(R_{h}-Z)_{x}\|_{H}^{2} + C_{\epsilon} h^{2(s-1)}\|Z\|_{W^{s,2}(S^{1})}^{2} \\
& \quad + C |t -\tau|
h^{s-1}\| z(\tau)\|_{W^{s,2}(S^{1})} ( \|(R_{h}-Z)_{x}\|_{H} + h^{s-1}\|Z\|_{W^{s,2}(S^{1})})
\end{align*}
where we have used \eqref{Ritz1}, \eqref{(4.2)extra} and \eqref{(4.2)bis} in the last inequality.
This yields
\begin{align*}
\|(R_{h}-Z)_{x}\|_{H} \leq C h^{s-1}\|Z\|_{W^{s,2}(S^{1})} + C |t -\tau|
h^{s-1}\| z(\tau)\|_{W^{s,2}(S^{1})} 
\end{align*}
which in turn gives \eqref{RH3}. Next we want to estimate the $L^{2}$-norm of $R_{h}-Z$ and ``win'' a factor $h$. To that end an Aubin-Nitsche trick must be employed.
With $ds:=ds(t)=|u_{x}(t)|dx$ and using Lax-Milgram we (weakly) solve the problem : find $w:S^{1} \to \R$ such that
\begin{align*}
\partial^{2}_{s} w = R_{h}-Z -\bar{m}, \qquad \int_{S^{1}}w ds=0,
\end{align*}
where (using the definition of the Ritz projection)
\begin{align*} \bar{m}:&=\int_{S^{1}}(R_{h}-Z) ds=\int_{S^{1}}(R_{h}-Z) |u_{x}(t)|dx= \int_{S^{1}} (z(\tau) -\mathcal{R}_{h}(z(\tau)))|u_{x}(t)|dx\\
&= \int_{S^{1}} (z(\tau) -\mathcal{R}_{h}(z(\tau)))(|u_{x}(t)|- |u_{x}(\tau)|)dx .
\end{align*}
Note that using the regularity of $u$ and \eqref{Ritz2} we have
\begin{align}\label{mbarest}
|\bar{m}| \leq C |t-\tau|\|z(\tau) -\mathcal{R}_{h}(z(\tau))\|_{H} \leq C |t-\tau| h^{s}\|z(\tau)\|_{W^{s,2}(S^{1})}.
\end{align}
Moreover, the Poincar\'{e}-inequality (on curves) gives $\| w\|_{L^{2}(ds(t))} \leq C\|\partial_{s} w\|_{L^{2}(ds(t))}$ (with a uniform constant $C$ independent of $t$ thanks to \eqref{regularity})
so that from the weak formulation we immediately infer
$\|\partial_{s} w\|_{L^{2}(ds(t))} \leq C \| R_{h}-Z -\bar{m} \|_{L^{2}(ds(t))} $.
Regularity theory finally yields
\begin{align*}
\| w \|_{W^{2,2}(ds(t))} \leq C \| R_{h}-Z -\bar{m} \|_{L^{2}(ds(t))}
\end{align*}
that is (using \eqref{regularity})
\begin{align}\label{reg-ww}
\| w \|_{W^{2,2}(S^{1})} \leq C \| R_{h}-Z -\bar{m} \|_{L^{2}(S^{1})}.
\end{align}
Exploiting the weak formulation for the solution $w$ we can write
\begin{align*}
\| R_{h}-Z -\bar{m}\|_{L^{2}(ds(t))}^{2} &=\int_{S^{1}} \partial_{s} w \partial_{s} (R_{h}-Z) ds
\\
&=\int_{S^{1}} \partial_{s} (w -w_{h})\partial_{s} (R_{h}-Z) ds + \int_{S^{1}} \partial_{s} w_{h}\partial_{s} (R_{h}-Z) ds=I+II
\end{align*}
where $w_{h} \in S_{h}$ will be chosen later. Using the definition of Ritz operator (and recalling that $ds=ds(t)$) we have that
\begin{align*}
II:&=\int_{S^{1}} \partial_{s} w_{h}\partial_{s} (R_{h}-Z) ds=\int_{S^{1}} \partial_{s} w_{h}\partial_{s} (z(\tau)-\mathcal{R}_{h}(z(\tau))) ds\\
&=\int_{S^{1}} w_{hx}\frac{(z(\tau)-\mathcal{R}_{h}(z(\tau)))_{x}}{|u_{x}(t)|} dx \\
&=\int_{S^{1}} w_{hx}\,(z(\tau)-\mathcal{R}_{h}(z(\tau)))_{x} \Big(\frac{1}{|u_{x}(t)|} - \frac{1}{|u_{x}(\tau)|} \Big) dx\\
&=\int_{S^{1}} (w_{hx}-w_{x})\,(z(\tau)-\mathcal{R}_{h}(z(\tau)))_{x} \Big(\frac{1}{|u_{x}(t)|} - \frac{1}{|u_{x}(\tau)|} \Big) dx \\
& \quad +
\int_{S^{1}} w_{x}\,(z(\tau)-\mathcal{R}_{h}(z(\tau)))_{x} \Big(\frac{1}{|u_{x}(t)|} - \frac{1}{|u_{x}(\tau)|} \Big) dx\\
& =
\int_{S^{1}} (w_{hx}-w_{x})\,(z(\tau)-\mathcal{R}_{h}(z(\tau)))_{x} \Big(\frac{1}{|u_{x}(t)|} - \frac{1}{|u_{x}(\tau)|} \Big) dx \\
&\quad -\int_{S^{1}} w_{xx}\,(z(\tau)-\mathcal{R}_{h}(z(\tau))) \Big(\frac{1}{|u_{x}(t)|} - \frac{1}{|u_{x}(\tau)|} \Big) dx \\
& \quad- \int_{S^{1}} w_{x}\,(z(\tau)-\mathcal{R}_{h}(z(\tau))) \Big(\frac{1}{|u_{x}(t)|} - \frac{1}{|u_{x}(\tau)|} \Big)_{x} dx.
\end{align*}
Therefore, using the regularity of $u$, \eqref{Ritz1}, \eqref{Ritz2}, we obtain
\begin{align*}
II &\leq C|\tau-t|h^{s-1}\| z(\tau)\|_{W^{s,2}(S^{1})} \|( w_{h} -w)_{x}\|_{L^{2}(S^{1})} \\
& \quad + C|\tau-t|h^{s}\| z(\tau)\|_{W^{s,2}(S^{1})} (\|w_{x}\|_{L^{2}(S^{1})} + \|w_{xx}\|_{L^{2}(S^{1})}).
\end{align*}
Choosing $w_{h}=I_{h}w$ and using \eqref{(4.2)}, \eqref{reg-ww}, \eqref{regularity}, we infer
\begin{align*}
II &\leq C|\tau-t|h^{s}\| z(\tau)\|_{W^{s,2}(S^{1})} \| R_{h}-Z -\bar{m} \|_{L^{2}(S^{1})},\\
I & \leq Ch\| R_{h}-Z -\bar{m} \|_{L^{2}(S^{1})} \| (R_{h}-Z)_{x}\|_{L^{2}(S^{1})}.
\end{align*}
Therefore
\begin{align*}
\| R_{h}-Z -\bar{m} \|_{L^{2}(S^{1})} \leq C|\tau-t|h^{s}\| z(\tau)\|_{W^{s,2}(S^{1})} + Ch
\| (R_{h}-Z)_{x}\|_{L^{2}(S^{1})}
\end{align*}
and then, using \eqref{mbarest},
\begin{align*}
\| R_{h}-Z \|_{L^{2}(S^{1})} \leq C|\tau-t|h^{s}\| z(\tau)\|_{W^{s,2}(S^{1})} + Ch
\| (R_{h}-Z)_{x}\|_{L^{2}(S^{1})}. 
\end{align*}
Together with \eqref{RH3} this yields \eqref{RH4}.
\end{proof}

\end{document}
